% Verbatim copy of /Users/alingalatan/workspaces/grothe/handoff/research-degree5/ai-notes/chatgpt-shared_driver_rectangle.tex (file dated 2026-11-05). Only this line was added.
% output/r789/shared_driver_rectangle.tex
\documentclass[11pt]{article}
\usepackage[T1]{fontenc}
\usepackage{lmodern,microtype,amsmath,amssymb,mathtools,booktabs}
\usepackage[a4paper,top=23mm,bottom=23mm,left=25mm,right=25mm,headheight=14pt]{geometry}
\usepackage{xcolor,fancyhdr}
\usepackage[colorlinks=true,linkcolor=ink,urlcolor=ink]{hyperref}
\definecolor{ink}{HTML}{233B53}\definecolor{muted}{HTML}{505C65}
\pagestyle{fancy}\fancyhf{}
\fancyhead[L]{\small\color{muted} Two shared-driver cubics on the remaining rectangle}
\fancyhead[R]{\small\color{muted} 5 November 2026}
\fancyfoot[C]{\small\thepage}\renewcommand{\headrulewidth}{0.3pt}
\setlength{\parindent}{0pt}\setlength{\parskip}{4pt}
\setlength{\abovedisplayskip}{5pt plus 2pt minus 1pt}
\setlength{\belowdisplayskip}{5pt plus 2pt minus 1pt}
\newcommand{\E}{\mathbb E}\newcommand{\R}{\mathbb R}\newcommand{\C}{\mathbb C}
\newcommand{\ind}{\mathbf1}
\newcommand{\head}[1]{\par\smallskip{\large\bfseries\color{ink}#1}\par\smallskip}
\newcommand{\subhead}[1]{\par\smallskip{\bfseries\color{ink}#1}\enspace}
\newcommand{\qed}{\hfill\(\square\)}
\hypersetup{pdftitle={Two shared-driver cubics on the remaining rectangle}}
\begin{document}
\thispagestyle{plain}
{\LARGE\bfseries\color{ink} Two shared-driver cubics\\[2pt] on the remaining rectangle\par}
\medskip
Put \(T=|Z|^2\), \(a=\ell-2\beta\), \(L=Z(a+\beta T)\), and
\[
 F(G)=\E|L+G|+\E|L-G|,\quad B(v)=2\E|L+\sqrt vW|,\quad
 \Lambda(G)=F(G)-\|G\|_2^2/\varepsilon.
\]
All complex standard Gaussians have \(\E|Z|^2=1\). The profile box is
\(.76\le\ell\le.90,\ .12\le\beta\le.16\), and \(.74\le\varepsilon\le.76\).
Let \(X,Y\) be independent remote complex Gaussians and define the unit-variance cubics
\[
 Q_1=X(|X|^2-2)/\sqrt2,\qquad Q_2=X(|Y|^2-1).
\]
For \(i=1,2\), let \(q=\sqrt u\,e^{i\phi}Q_i\). The remote linear form \(U\) may
use any remote coordinates, with any coefficients, and \(w=\E|U|^2\).

\head{Results}
On the entire rectangular range
\[
 \boxed{\qquad .00551\le u\le.0463,\qquad 0\le w\le.0223,\qquad}
\]
including every relative phase, we prove the stronger same-variance statement
\[
 \boxed{\quad F(q+(T-1)U)\le B(u+w)-10^{-6}\ell.\quad} \tag{1}
\]
Thus these two shared-driver families satisfy exact K on the remaining rectangle.
The proof includes exact integer enclosures, replayed by the accompanying
\texttt{verify\_shared.mjs}. It is independent of the remote-only \(10^{-5}\) input.

The corresponding content inequality also holds uniformly on the profile box:
\[
 \boxed{\quad
 \|\Pi_{1,0}F_A(q)\|_2^2
 \le \frac{19}{100}\bigl(B(u)-F(q)\bigr)-3\cdot10^{-6},
 \qquad F_A(q)=\E_Z[(T-1)\operatorname{ph}(L+q)] .
 \quad} \tag{2}
\]
The right side uses the same-variance remote deficit. Since
\(B(u)-F(q)\le S-\Lambda(q)\), (2) implies the requested content inequality.

At the three application rows we additionally extend the zero-buffer A cost to
\(\|U\|_2\le.15\), for \emph{any} remote \(q\) and any dependence with \(U\):
\[
 \Lambda(q+(T-1)U)\le\Lambda(q)+2\Re\langle U,F_A(q)\rangle-\tfrac{19}{100}w.
 \tag{3}
\]
Completing the square with (2) yields \(\Lambda(q+(T-1)U)\le S-3/190000\) at those rows.

For R7 we give an explicit sixth-order error and its radius on pages 2--3.
For R8, page 4 identifies the extra term in the shared variance derivative and gives
a direct conditional-complement estimate. Arbitrary shared cubic tensors beyond
the stated families remain open in the target rectangle.

\newpage
\head{R7. A uniform sixth-order remainder}
This estimate needs only finite sixth moments of the remote pair. Write
\(u=\E|q|^2,\ w=\E|U|^2,\ v=u+w\), and suppose \(\E[\bar q U]=0\).
Let \(C_2,C_4\) be the matched coefficients from the preceding note, so
\[
 C_2=\frac w2\int_0^\infty\frac{t(t-2)e^{-t}}{\sqrt t(a+\beta t)}\,dt,\qquad
 C_4=\frac1{32}\mathcal J[\E|q+(t-1)U|^4-2v^2].
\]
Here \(\mathcal J[P]=\int_0^\infty[e^{-t}P(t)/(a+\beta t)^3-P(0)/a^3]t^{-3/2}\,dt\).
Define
\[
 R=\frac{|q|}{.44}+\frac{|U|}{.12},\qquad
 D_6=7\E R^6+\frac{42v^3}{.44^6}.
\]
Then, uniformly on the box and for every real \(\lambda\),
\[
 \boxed{\quad
 \left|F(\lambda[q+(T-1)U])-B(\lambda^2v)
       -C_2\lambda^2-C_4\lambda^4\right|
 \le D_6|\lambda|^6 .\quad} \tag{4}
\]
Using Minkowski, the computable upper bound
\[
 D_6\le
 7\left[
 \left(\frac{(\E|q|^6)^{1/6}}{.44}
        +\frac{6^{1/6}\sqrt w}{.12}\right)^6
 +\frac{6v^3}{.44^6}\right] \tag{5}
\]
uses only individual sixth moments. Fixed-chaos hypercontractivity gives
\((\E|q|^6)^{1/6}\le5^{3/2}\sqrt u\) for a cubic; sharper available moments can replace it.

\subhead{Proof.}
Average real projections using \(|z|=(\pi/2)\E_\theta|\Re(e^{-i\theta}z)|\).
Write \(Z=x+iy\), where \(x,y\) have density \(e^{-x^2}/\sqrt\pi\), and put
\(b=\Re(e^{-i\theta}q), d=\Re(e^{-i\theta}U)\),
\[
 A=a+\beta y^2,\quad \eta=\lambda[b+d(y^2-1)],\quad k=\lambda d,\quad
 x_0=|\eta|/A,\quad z=|\lambda|R.
\]
The gain in \(F\) equals \(2\sqrt\pi\) times the expectation in \(y,\theta,\) and
the remote variables of
\(\int_0^\infty e^{-x^2}(|\eta+kx^2|-Ax-\beta x^3)_+\,dx\).
We have \(x_0\le z,\ |k|\le .12z\).

If \(k^2\le A\beta\), only the branch with the sign of \(\eta\) can contribute,
and its support is contained in \([0,4x_0/3]\). Put
\(\widetilde k=k\,\operatorname{sgn}\eta\).
Linearizing the positive part about \((|\eta|-Ax)_+\), then replacing
\(e^{-x^2}\) by \(1-x^2\) in the base term and by 1 in the correction, gives
\[
 P=\frac{Ax_0^2}{2}+\frac{\widetilde kx_0^3}{3}
                  -\frac{\beta x_0^4}{4}-\frac{Ax_0^4}{12}.
\]
For \(M=(16/9)|k|x_0^2+(64/27)\beta x_0^3\le(16/27)z^3\), the error is at most
\[
 2M^2/A+Ax_0^6/60+|k|x_0^5/5+\beta x_0^6/6
 <(5/3+A/60)z^6.
\]
The positive-part error is supported on an interval of length at most \(2M/A\)
and bounded by \(M\). The other terms use \(0\le e^{-t}-1+t\le t^2/2\).

\newpage
\head{R7. Radius and what it covers}
If \(k^2>A\beta\), then \(z>\sqrt{.44/.12}>1.9\). Bound the actual positive-part
integral by \((\sqrt\pi/2)(Az+|k|/2)\) and bound the four terms of \(P\) separately.
After division by \(z^6\), this gives an error below \(1/30+A/10\).
Combining the two cases, averaging
\(\E_y A=a+\beta/2\le.74\), and using \(\sqrt\pi<1.8\), bounds the error in
\(F\) by \(7\E R^6|\lambda|^6\). The polynomial contains only degrees two and four;
uniqueness of the established expansion identifies its coefficients.
Applying the same estimate to the independent Gaussian comparator, whose sixth
moment is \(6v^3\), proves (4).\qed

For \(w>0\), \(C_2<0\). The explicit sufficient radius from (4) is
\[
 0<|\lambda|<\lambda_{\rm cert},\qquad
 \lambda_{\rm cert}^2=
 \frac{-C_4+\sqrt{C_4^2-4D_6C_2}}{2D_6}. \tag{6}
\]
When \(w=0,\ C_4<0\), use \(\lambda_{\rm cert}^2=-C_4/D_6\).
The bound \(D_6\) from (5) may be substituted in (6).

The radius does not reach the target rectangle for either named example.
Here is an analytic check, without relying on a numerical scan. Put
\(z_6=(\E R^6)^{1/6}\). On the box
\[
 |C_2|<.03z_6^2,\qquad |C_4|<z_6^4 .
 \tag{7}
\]
For the first estimate, use \(|C_2|\le w\sqrt\pi/(2\cdot.44)\) and
\(w\le.12^2z_6^2\). For the second, the formula for \(\mathcal J[(t-1)^j]\)
in the preceding note gives its absolute value below 275 for \(0\le j\le4\).
Indeed its normalized polynomial is bounded by 9, and
\(\E\sqrt{1+(\beta/a)Y}<1.45\) for \(Y\sim\Gamma(3,1)\).
The sum of absolute fourth-moment coefficients is at most
\(\E(|q|+|U|)^4\), and \(v^2\le.44^4 z_6^4\). Thus the quartic bound is at most
\((825/32).44^4z_6^4<z_6^4\).

The normalized sixth moments of \(Q_1,Q_2\) are respectively \(11352,1590\),
by exponential moments. Since \(u\ge.00551\), either family has \(z_6>1/2\).
Therefore \(D_6\ge7z_6^6>|C_2|+|C_4|\) at every point of the rectangle.
The sufficient inequality from (4) fails at \(\lambda=1\).

As a scale illustration, take a normalized direction \(u=.8,w=.2\) at the peak
row, with aligned positive coefficients. Ordinary evaluation of the formulas gives
\(\lambda_{\rm cert}\simeq .00594\) for \(Q_1\) and \(.00805\) for \(Q_2\), using (5).
The point \(u_q=.02,w=.005\) in the requested rectangle has common amplitude
\(\sqrt{.025}\simeq.1581\). The proof of (1) below avoids this Taylor loss.

\newpage
\head{R8. Shared differentiation and conditional complements}
Take \(U=\sqrt w X\) after a unitary change of remote coordinates. A conditional
mean-zero complement \(J=q-\E[q\mid X]\) can depend on \(X\).
Even if \(J\) is independent of \(X\), the component \(q_0(X)=\E[q\mid X]\)
still contributes to variance differentiation.

For a smooth function \(\Phi:\R^2\to\R\), write \(h=T-1\) and
\(Y_v=L+q(X,\text{other coordinates})+\sqrt v\,hX\), identifying complex numbers
with real two-vectors.
\href{https://arxiv.org/abs/0712.2940}{Gaussian integration by parts} with
\({\rm Cov}(X)=I/2\) gives the exact formula
\[
 \frac d{dv}\E\Phi(Y_v)
 =\frac14\E[h^2\Delta\Phi(Y_v)]
 +\frac1{4\sqrt v}\E\!\left[h
       \sum_{i,j}(\partial_jq_i)(\partial_{ij}\Phi)(Y_v)\right]. \tag{8}
\]
A fresh Gaussian smoothing of the norm makes this identity directly applicable.
The last term vanishes for the disjoint variable in the prior proof. It remains
for \(q_0(X)\), and a conditional mean-zero \(J\) need not cancel its own
derivatives against the Hessian depending on \(J\).

There is a cubic obstruction even with \(J=0\). For \(q=cX(|X|^2-2)\) and small
real \(c>0\), the expansion of the signed visible potential gives
\[
 \E[\bar X F_A(q)]
 =2\mathcal J[t-1]c^3+O(c^5)>0 .
\]
Here \(\E[\bar X H|H|^2]=32\) for \(H=X(|X|^2-2)\), and
\(\mathcal J[t-1]>0\). Thus adding a sufficiently small positive
\((T-1)\alpha X\) increases \(\Lambda\) to first order in \(\alpha\).
This prevents the independent-shell decrease from carrying over through the
proposed split alone.

\subhead{A direct zero-buffer transfer.}
There is nevertheless a quantitative estimate for the complement. Let
\(q_0,U\) be measurable with respect to the conditioning data,
\(\E[J\mid\text{data}]=0\), \(\delta=\E|J|^2\), and \(w=\E|U|^2\).
For any \(0<k_0<\sqrt{3a\beta}\), put \(d_0=a-k_0^2/(3\beta)\). Then
\[
 \begin{aligned}
 \Lambda(q_0+J+(T-1)U)
 \le{}&\Lambda(q_0+(T-1)U)
       +\left(\frac{\sqrt\pi}{2d_0}-\frac1\varepsilon\right)\delta\\
 &+\sqrt2\,\pi\sqrt\delta\,
       \sqrt{\operatorname{erfc}(k_0/\sqrt w)} .
 \end{aligned} \tag{9}
\]
At \(w=0\) interpret the tail as zero. To prove it, use the real projection in
the next section. When \(|\Re U|\le k_0\), the visible cubic in \(x\) has derivative
at least \(d_0+\beta y^2\), so its density is bounded by
\(1/(\sqrt\pi d_0)\). Taylor in the constant shift \(\Re J\) then bounds the
remainder by this density bound times \((\Re J)^2\). On the exceptional event
use \(2|\Re J|\). Average projections and apply Cauchy--Schwarz; the Gaussian
tail probability is \(\operatorname{erfc}(k_0/\sqrt w)\).
The linear term cancels by the conditional mean-zero hypothesis.
Estimate (9) gives a zero-buffer alternative, while a general large complement
still needs additional structure.

\newpage
\head{A projected quadratic cost for the entire target amplitude range}
For \(q\) arbitrary remote and a linear Gaussian \(U\), use real projections again.
Put \(\eta=\Re q,\ k=\Re U\), after the projection rotation, and
\[
 D=x(a+\beta(x^2+y^2)),\qquad
 P_{a,\beta}(z)=\int_0^\infty
   \sup_{x\ge0}\frac{(x^2+y^2-1)^2e^{-x^2-y^2}}
                     {a+\beta y^2+3\beta x^2-2zx}\,dy .
\]
The denominator is positive for \(z<\sqrt{3a\beta}\).
For \(|k|\le k_0\), differentiation of the real absolute value gives
\[
 \begin{aligned}
 &F(q+(T-1)U)-F(q)-2\Re\langle U,F_A(q)\rangle\\
 &\quad\le 4\E\left[k^2\int_0^1(1-s)P_{a,\beta}(s|k|)\,ds\,
                   \ind_{\{|k|\le k_0\}}\right]
 +\frac{4\sqrt\pi}{e}\sqrt w\,e^{-k_0^2/w}.
 \end{aligned}\tag{10}
\]
For each \(y,\eta\), the equation
\(D+sk(x^2+y^2-1)+\eta=0\) has a unique simple root.
Its weighted density is bounded by \(2P_{a,\beta}(s|k|)/\pi\).
The second derivative of the real absolute mean is twice that density times \(k^2\),
and the factor \(\pi\) from real projections gives (10).
For the exceptional event, the first-order remainder is at most
\(2|k|\E|T-1|=4|k|/e\) before multiplying by \(\pi\).
Use \(\E[|k|\ind_{\{|k|>k_0\}}]=\sqrt{w/\pi}\,e^{-k_0^2/w}\).

The envelope \(P_{a,\beta}\) is convex in \(z\). If
\(P_{a,\beta}(z)\le A+Bz+Cz^2\) on \([0,k_0]\), then (10), divided by \(w\), is at most
\[
 A+\frac{2B\sqrt w}{3\sqrt\pi}+\frac{Cw}{4}
       +\frac{4\sqrt\pi}{e\sqrt w}e^{-k_0^2/w}. \tag{11}
\]
At the smallest of the three actual profiles,
\((a,\beta)=(.5592,.1432)\), the integer checker verifies
\[
 P(0)<1.07,\qquad P(.2)<1.18,\qquad P(.42)<1.74 .
\]
The two chords are bounded above by \(1.07+.55z+3z^2\).
For \(\sqrt w\le.15\), (11) is smaller than
\(1/.7564-.19\). The other two rows increase both \(a,\beta\) and have a
penalty at least \(1/.7564\). This proves (3) with zero additive error.

\subhead{Normalized profiles.}
For the full box put \(p=\beta/\ell\), \(a/\ell=1-2p\), and
\(\rho=\sqrt w/\ell\). On \(2/15\le p\le4/19,\ \rho\le.197\), the same argument
gives, for arbitrary sharing,
\[
 \frac{F(q+(T-1)U)-F(q)-2\Re\langle U,F_A(q)\rangle}{\ell}
 \le K(p)\rho^2,\qquad
 K(p)=.940+.105\,\frac{p-2/15}{4/19-2/15}. \tag{12}
\]
The constants and their verification follow on the next page.

\newpage
\head{Verification of the projected costs}
For \(P_{1-2p,p}\) at the two extreme values of \(p\), the checker verifies the
following upper bounds, each rounded upward.
\[
 \begin{array}{c|rrrr|rrr}
 \toprule
 p & P(0)&P(.2)&P(.4)&P(.5)&A&B&C\\ \midrule
 2/15&.856&.925&1.170&1.750&.856&.400&3.300\\
 4/19&.982&1.044&1.220&1.430&.982&.320&1.600\\ \bottomrule
 \end{array}
\]
On each interval between the listed \(z\)'s, convexity bounds \(P\) by the chord.
The displayed polynomial \(A+Bz+Cz^2\) dominates every such chord.
The checker verifies that quadratic inequality at its endpoints and, when
appropriate, at its vertex.

At \(\rho=.197\), the last term of (11) with \(k_0=.5\) is below .022.
The function \(\rho^{-1}e^{-k_0^2/\rho^2}\) increases on the range.
Use \(\sqrt\pi>1.77\) in the cubic-moment term. The upper bounds in (11) are
then respectively .940 and 1.045. The envelope is convex in \(p\) as well:
the positive denominator of each integrand is affine in \(p\), and integration
and suprema preserve convexity. Interpolation in \(p\) proves (12).

\subhead{Enclosing \(P\).}
The supremum integral is enclosed on \(0\le x,y\le4\) with mesh \(1/256\).
On a rectangle use the larger endpoint of \((x^2+y^2-1)^2\), the lower endpoints
in the decreasing exponential, and
\(a+\beta y_{\min}^2+3\beta x_{\min}^2-2k x_{\max}\) in the denominator.
For each \(y\)-interval take the largest bound among the \(x\)-intervals.
The denominator lower bound is positive, as checked.
For \(x\ge4\) the numerator decreases and the denominator increases, so the
boundary \(x=4\) covers this tail. For \(y\ge4\) the denominator is at least 1,
and the supremum of the numerator is at most \(y^4e^{-y^2}\). Its integral is
below \(10^{-5}\), the amount added by the checker.

\head{R9. Exact reductions for the two remote cubics}
Define the circular mean \(H(s)=\E_\theta|1+se^{i\theta}|\), and
\[
 j(s)=\frac{H(s)-1}{s},\qquad h(s)=H'(s),\qquad j(0)=h(0)=0 .
\]
Both increase from 0 to 1, by convexity of \(H\).
Write \(R_p(t)=\sqrt t[1+p(t-2)]\) and \(\nu=\sqrt u/\ell\).
For the two templates put
\[
 B_1=|X|\,\lvert |X|^2-2\rvert/\sqrt2,\quad
 B_2=|X|\,\lvert |Y|^2-1\rvert,\quad
 s_1=\operatorname{sgn}(|X|^2-2),\quad s_2=\operatorname{sgn}(|Y|^2-1).
\]
Then the normalized remote mean gain and the content coefficient are exactly
\[
 \frac{F(q)-F(0)}{2\ell}
   =\nu\,\E[B_i j(\nu B_i/R_p(T))],\qquad
 \xi_i=\E[s_i|X|(T-1)h(\nu B_i/R_p(T))].
 \tag{13}
\]
The projection has only an \(X\)-coefficient, by rotating each remote coordinate,
so \(\|\Pi_{1,0}F_A(q)\|_2=|\xi_i|\). The phase \(e^{i\phi}\) rotates this coefficient.

\newpage
\head{Positive enclosures for the content and deficit}
The Gaussian comparator uses \(B_{\rm G}=|W|\) in the first expression in (13).
For a lower bound on its gain and an upper bound on the cubic gain, the remaining
integrands are nonnegative and increasing in the ratio and in the factor \(B_i\).
For the content, separate the four positive integrals according to the signs
of \(s_i\) and \(T-1\), taking the difference only after enclosing each integral.
This is the point at which the cancellations required for a small first-chaos
projection are retained.

Let \(b=1025/1024\). Radial coordinates \(|Z|,|X|,|Y|\) are divided into
mesh-\(1/1024\) intervals on \([0,5]\).
Their masses are differences of \(e^{-x^2}\), and visible signed masses use
the endpoint antiderivative \(-t e^{-t}\).
Use endpoint bounds for factors \(|X|\), \(\lvert |X|^2-2\rvert\),
\(\lvert |Y|^2-1\rvert\), and \(R_p(T)\); endpoints in \(p\) enclose a full profile interval.
If a cell crosses the sign change for \(s_1\), assign its full upper mass to
both sign components, and assign zero to both lower components.
Push the bounds on all positive radii to enclosing powers \(b^j\).

For each resulting positive integral a missing-mass budget of \(10^{-5}\) suffices.
Indeed a remote tail event has probability at most \(2e^{-25}\), and
\(\E B_i^2=\E|X|^2=1\). Cauchy--Schwarz bounds its contribution by
\(\sqrt{2e^{-25}}<5.3\cdot10^{-6}\).
The first visible cell, whose lower radius is zero, has mass below \(1024^{-2}<10^{-6}\).
The visible outer tail has mass below \(e^{-25}\), or weighted mass
\(25e^{-25}\) for the content. The total is below \(10^{-5}\).
In the signed content, one such budget bounds the total omitted variation.

\subhead{The angular tables.}
For \(0\le z\le1\), set
\[
 A(z)=\sum_{k\ge1}a_kz^{2k-2},\quad C(z)=\sum_{k\ge1}2ka_kz^{2k-2},\quad
 D(z)=\sum_{k\ge1}(2k-1)a_kz^{2k-2},
\]
where \(a_k=\binom{1/2}{k}^2\). Thus \(a_1=1/4\) and
\(a_{k+1}/a_k=(2k-1)^2/(2k+2)^2\). If \(z\le1\),
\[
 j(z)=zA(z),\quad h(z)=zC(z),\qquad
 j(1/z)=1-z+z^2A(z),\quad h(1/z)=1-z^2D(z).
\]
Every coefficient is positive. The exact sums at 1 are
\(A(1)=4/\pi-1,\ C(1)=2/\pi,\ D(1)=1-2/\pi\), obtained directly by integrating
\(2|\cos(\theta/2)|\).
Truncate at 64 terms for a lower bound. For the upper bound, add \(z^{128}\)
times an upper bound on the omitted sum at 1. The checker uses respectively
\(.273240,.636620,.363381\) as upper bounds on the full sums.
This encloses \(j,h\) at every required logarithmic-grid argument, including near 1.

\newpage
\head{Exact arithmetic and the content certificate}
All numerical bounds in the checker are integers over \(2^{56}\).
For \(e^{-x}\), \(0\le x\le25\), scale by 64, bound by alternating Taylor
partial sums of orders 13 and 14, then square the interval six times.
The powers \(b^j\) are enclosed by directed multiplication or division.
Approximate logarithms and square roots supply starting guesses only; the
returned indices and integer square roots are determined by exact comparisons.

Positive histograms on the logarithmic grid combine by convolution.
To perform these convolutions exactly, the code packs each integer coefficient
into a 128-bit digit of a single integer and multiplies the two integers.
Before every multiplication it checks
\[
 \min\{\|A\|_1\|B\|_\infty,\ \|A\|_\infty\|B\|_1\}<2^{128},
\]
so no coefficient can carry into its neighbor. It then extracts the digits and
rounds the common scaling denominator outward. This is exact polynomial
convolution using ordinary integer multiplication.

On each band \(b^n\le\nu\le b^{n+1}\), use the lower endpoint for each lower
positive integral and the upper endpoint for each upper positive integral.
For the content subtract the appropriate upper bounds on negative components
from lower bounds on positive components, and conversely. Let \(\xi^*\) bound
\(|\xi_i|\), let \(G_{\rm G}^{-}\) be the lower Gaussian gain factor, and
\(G_i^{+}\) the upper cubic gain factor in (13), including the tail budget.
It suffices to verify
\[
 (\xi^*)^2 <
 .19\cdot2(.76)b^n(G_{\rm G}^{-}-G_i^{+}) .
\]
The factor .76 handles every \(\ell\) in the box.
The replay covers 20 profile intervals and each of the 1351 amplitude bands
with \(-2600\le n\le-1250\). This enlarged range contains
\(\sqrt u/\ell\) throughout the stated rectangle, checked by exact squared inequalities.
The smallest positive enclosure, in units of \(2^{-56}\), is
\[
 241420829874/72057594037927936>3\cdot10^{-6}.
\]
This proves (2). It also proves a strictly positive same-variance remote deficit
for these examples on the range.\qed

\subhead{General tensors.}
For a general Wick cubic, Gaussian integration by parts gives the exact first-chaos
coefficient \(\E[\bar z_jF_A(q)]=\E[\partial_{z_j}(F_A(q))]\).
If \(F_A(y)=y\,g(|y|^2)\), this is
\[
 \E\big[(g(s)+sg'(s))\partial_{z_j}q
                 +q^2g'(s)\partial_{z_j}\bar q\big],\qquad s=|q|^2.
\]
The constant part cancels because the derivatives of \(q\) are centered quadratics.
The identity exposes the cubic-order content, but a sharp bound comparing this
expression to the remote deficit for arbitrary tensors remains unresolved here.
The certificate above uses the specified templates and their torus symmetries.

\newpage
\head{Same-variance domination for the complete shared perturbation}
We now combine the exact content enclosures with (12) directly.
Let \(\rho=\sqrt w/\ell\). On the target rectangle \(\rho<.197\).
For a profile interval and a band \(b^n\le\nu\le b^{n+1}\), let \(c^*\) be the
upper bound on \(|\xi_i|\), \(g_i^+\) the upper bound on the remote gain factor
including the missing mass, and \(K^+\) the upper endpoint of \(K(p)\).
For a shell-radius interval \(\rho_0\le\rho\le\rho_1\), the upper bound for the
normalized gain over \(F(0)\) is
\[
 2b^{n+1}g_i^++2\rho_1c^*+K^+\rho_1^2 .
 \tag{14}
\]
By Cauchy--Schwarz, \(2\Re\langle U,F_A(q)\rangle/\ell\le2\rho|\xi_i|\),
so this holds for every phase and every remote linear form \(U\).

Choose a lower grid point \(b^m\le\sqrt{b^{2n}+\rho_0^2}\).
The normalized Gaussian gain at the matched variance is bounded below by
\(2b^m G_{\rm G}^{-}(b^m)\), using monotonicity in the Gaussian amplitude.
The checker verifies that this lower bound minus (14) exceeds \(10^{-6}\).
It uses \(-2560\le n\le-1280\), and shell endpoints on the mesh \(1/2048\)
with final endpoint .197. The two end-amplitude tests in the script verify
that these bands contain every \(\sqrt u/\ell\) in the target rectangle.
Shell intervals are bisected until the inequality is certified or a mesh cell is
reached, with further profile subdivision when necessary.

The smallest certified gap in the replay is
\[
 \frac{72059660296}{72057594037927936}>10^{-6}.
\]
Multiplying by \(\ell\) proves (1).\qed

\head{Replay and scope}
Run \texttt{node verify\_shared.mjs} using Node.js 20 or later.
The script has no external dependencies. Its final checks also establish the
projected-envelope bounds, the Gaussian tail constants, and (3).

The proven range for the two named remote templates is the full rectangle on
page 1, larger than the surviving sloping region in the question.
For a general shared cubic and linear form, we have the explicit Taylor estimate,
the corrected shared differentiation identity, the zero-buffer complement transfer,
and the strengthened A cost at the three rows. The general content inequality
and general all-amplitudes K on the surviving rectangle remain open here.

The previous disjoint theorem uses the project's remote-only input. With its
certified error \(10^{-5}\), the resulting disjoint conclusion has error at most
\(10^{-5}\). The explicit shared-template comparison (1) is exact.
\end{document}
